\begin{frame}[t]{A drug combination preserved cell killing while reducing PD-L1 induction.\ev{\tagO}}
\vspace{0pt}
\lead{Measure cell viability and PD-L1, an immune-checkpoint ligand.}
\begin{center}\begin{tikzpicture}[x=1cm,y=1cm,every node/.style={align=center,font=\fontsize{12.5}{15}\selectfont}]
\node[state,text width=5.3cm,minimum height=1.8cm] (single) at (0,0) {\textbf{Single-drug response}\\[5pt]PD-L1 mRNA rose about\\\textbf{10--100 times} in E0771 cells};
\node[evidence,text width=5.3cm,minimum height=1.8cm] (combo) at (6.6,0) {\textbf{Abemaciclib + SI-2}\\[5pt]Lower PD-L1 induction\\with cytotoxicity preserved};
\draw[flow](single)--(combo);
\node[text width=12.5cm,font=\fontsize{11.5}{14}\selectfont,text=Muted] at (3.3,-1.8) {E0771 mouse breast-cancer cells; 72-hour combination assay.};
\node[text width=12.5cm,font=\fontsize{13}{16}\selectfont,text=Teal] at (3.3,-2.55) {Cell killing and immune-related expression are separate objectives.};
\end{tikzpicture}\end{center}
\sources{Gilad, Eliaz et al., \href{https://doi.org/10.1038/s41598-019-51537-7}{Scientific Reports 9:15099 (2019)}, Figs. 1, 3 and S2.}
\qadetail{Primary full text: https://pmc.ncbi.nlm.nih.gov/articles/PMC6805932/. In four human/mouse breast-cancer lines, drug exposure near IC50 induced PD-L1 expression. In E0771 cells, PD-L1 mRNA induction varied by drug, approximately tenfold with paclitaxel or abemaciclib and up to approximately one hundredfold with doxorubicin or topotecan. This range concerns different single-drug exposures, not the fold reduction from the combination. After 72 hours in E0771 cells, abemaciclib with SI-2 reduced PD-L1 mRNA/surface protein induction toward the weaker-inducer level while retaining cytotoxicity. The text gives no exact combination fold reduction, and the slide invents none. Addition of abemaciclib up to 625 nM did not further elevate ATF4 stress protein; retained as a separate endpoint. The March 2020 correction changes the Figure 2A label INF to IFN, not these results. These are tumor-cell marker and viability assays, not measurements of T-cell response, patient benefit, a deployed treatment policy or autonomous immunotherapy. Eliaz is a coauthor; Gilad performed experiments and Eliaz contributed analysis, visualization and study conception. The result motivates the user's response-guided immunotherapy research direction by showing why multiple measured objectives matter. It is distinct from the later SI-12 synthetic-lethal CRISPR screen.}
\note{[0:45] The connection to immunotherapy starts with a measured design tradeoff. PD-L1 is an immune-checkpoint ligand. In E0771 mouse breast-cancer cells, different drugs increased its messenger RNA roughly tenfold to a hundredfold. Combining abemaciclib with SI-2 lowered that induction while preserving cell killing in a seventy-two-hour assay. We measured cancer-cell viability and expression. The research direction is to guide future combinations from observed responses. For computation, the lesson is similar: evaluate both the intended outcome and the changes your action induces.}
\end{frame}
